% Brouillon de la partie finale « Pour aller plus loin » (17/11/2026). Chiffres : sorties/pal.json
% (calculs/pal_calculs.py) ; compte rendu : sorties/pal.md ; annexe : annexes/annexe_pal.tex. Hors scénario central.
% Figure facultative : figures/gpal_revenus_irreguliers (macro \graphique du préambule).
\section{Pour aller plus loin : récompenser ceux qui construisent et ceux qui prennent des risques}

Trois prolongements ne font pas partie de notre scénario chiffré. Ils ont un coût, que nous indiquons ; aucun n'est
financé par une croissance espérée.

\textbf{Notre réforme doit inciter non seulement les propriétaires à construire, mais aussi les maires à les laisser
construire.} Depuis la suppression de la taxe d'habitation, un logement neuf ne rapporte à sa commune que la taxe
foncière, après deux ans d'exonération ; la Cour des comptes (2025) juge que cela a réduit l'incitation à délivrer des
permis. Notre réforme aggrave ce point : en taxant le sol quatre fois plus que le bâti, elle fait passer la recette
annuelle qu'un appartement neuf de 290\,000~€ apporte au bloc communal d'environ 780~€ à 310~€. En revanche, la commune
gagnera désormais quand elle ouvre des droits à construire, puisque la valeur du sol montera : c'est la décision qui
compte. Pour compenser le reste, nous proposons une dotation automatique de 5\,000~€ par logement supplémentaire, net
des démolitions, dans les zones tendues, versée l'année où il entre au rôle de la taxe foncière, sans autre condition.
Pour 150\,000 à 300\,000 logements par an, elle coûterait 0,75 à 1,5~Md€, pris sur la dotation globale de
fonctionnement des communes à enveloppe constante : les communes qui ne construisent pas y perdraient. L'Angleterre a
pratiqué une prime de ce type depuis 2011 ; son évaluation n'a pas démontré d'effet sur la construction. Nous
n'en chiffrons donc aucun.

\textbf{L'État partage immédiatement les bénéfices des entreprises, mais beaucoup moins facilement leurs pertes.} Un
bénéfice est taxé l'année même ; une perte ne réduit l'impôt que plus tard, sans intérêts, dans la limite de 1~M€ plus
la moitié du bénéfice au-delà, ou jamais si l'entreprise disparaît. Pour une jeune entreprise qui perd de l'argent trois
ans avant d'en gagner, le taux effectif de l'impôt sur son projet passe ainsi de 25 à près de 27~\%, et notre déduction
immédiate de 50~\% perd les trois quarts de sa valeur. Le remède le plus simple consiste à revaloriser chaque année les
nouveaux déficits reportables au taux auquel l'État emprunte. Pour une entreprise qui finit par être bénéficiaire, le
délai ne coûte alors plus rien ; à celle qui échoue, l'État ne verse rien, ce qui évite de maintenir en vie des
entreprises non viables, risque que la littérature a documenté pour les remboursements immédiats de pertes. Ce n'est
pas gratuit : environ 1~Md€ par an une fois la mesure montée en charge (0,4 à 2,2 selon des données que
l'administration ne publie pas), presque rien les premières années. Les études montrent que les entreprises prennent
plus de risques quand l'État partage leurs pertes ; elles ne permettent pas de chiffrer un gain pour la France.

\textbf{À revenu cumulé identique, notre fiscalité pénalise celui qui prend le risque d'avoir des revenus irréguliers.}
Avec le nouveau barème, Antoine, qui gagne 50\,000~€ par an pendant cinq ans, paie 48\,700~€ d'impôt. Béatrice, qui
lance son activité et gagne successivement 10\,000, 20\,000, 30\,000, 70\,000 puis 120\,000~€, soit le même total, paie
60\,000~€ : 11\,300~€ de plus, parce qu'un impôt progressif taxe davantage les bonnes années qu'il n'allège les
mauvaises. Le droit actuel ne corrige ce défaut que pour des cas particuliers : agriculteurs, artistes, écrivains,
sportifs, revenus exceptionnels hors activité normale. Nous proposons de remplacer ces régimes par une option ouverte à
tous les revenus professionnels : pendant cinq ans, chacun paie l'impôt qu'il aurait dû sur son revenu moyen,
régularisé chaque année et à la sortie (cessation, décès, départ à l'étranger), le prélèvement à la source suivant ce
revenu moyen. L'écart entre Antoine et Béatrice disparaît (figure~\ref{fig:gpal}) ; un lissage sur trois ans le réduirait
d'un tiers. Ce que l'État y perd, la surtaxe de l'irrégularité, ne peut être chiffré sans suivre les mêmes contribuables
plusieurs années ; les régimes actuels de lissage coûtent peu (11~M€ pour celui des agriculteurs).

\graphique{gpal_revenus_irreguliers}{À revenu cumulé égal, des revenus irréguliers coûtent 11\,000~€ d'impôt de plus en
cinq ans ; une option de lissage sur cinq ans efface l'écart}{fig:gpal}
